\documentclass[11pt,a4paper]{report}
\usepackage[utf8]{inputenc}
\usepackage[T1]{fontenc}
\usepackage[italian]{babel}
\usepackage{lmodern}
\usepackage[margin=1.8cm]{geometry}
\usepackage{amsmath,amssymb,amsthm,mathtools}
\usepackage{booktabs}
\usepackage{tabularx}
\usepackage{xcolor}
\usepackage{enumitem}
\usepackage[most]{tcolorbox}
\usepackage{hyperref}

\hypersetup{
    colorlinks=true,
    linkcolor=blue!70!black,
    citecolor=blue!70!black,
    urlcolor=blue!70!black
}

% Palette Colori
\definecolor{primaryblue}{RGB}{24, 76, 120}
\definecolor{boxbg}{RGB}{248, 250, 252}
\definecolor{boxframe}{RGB}{52, 116, 170}
\definecolor{intbg}{RGB}{240, 253, 244}
\definecolor{intframe}{RGB}{34, 139, 34}
\definecolor{usebg}{RGB}{238, 246, 255}
\definecolor{useframe}{RGB}{41, 128, 185}
\definecolor{demobg}{RGB}{250, 245, 255}
\definecolor{demoframe}{RGB}{128, 0, 128}
\definecolor{oralbg}{RGB}{255, 251, 235}
\definecolor{oralframe}{RGB}{217, 119, 6}

% Box Teorema Principale
\newtcolorbox{teoremabox}[2][]{
    enhanced, breakable,
    colback=boxbg, colframe=boxframe,
    fonttitle=\bfseries\color{white},
    title={#2}, arc=2mm, boxrule=1.2pt,
    left=3.5mm, right=3.5mm, top=3mm, bottom=3mm,
    before skip=10pt, after skip=6pt, #1
}

% Sotto-Box 1: Cosa fa (Intuizione)
\newtcolorbox{cosafa}[1][Cosa Fa (L'Intuizione in parole semplici)]{
    enhanced, breakable,
    colback=intbg, colframe=intframe,
    fonttitle=\bfseries\small\color{white},
    title={$\blacktriangleright$ #1}, arc=1.5mm, boxrule=0.8pt,
    left=2.5mm, right=2.5mm, top=2mm, bottom=2mm,
    before skip=4pt, after skip=4pt
}

% Sotto-Box 2: A cosa serve (Utilizzo Pratico & Progetto SIR)
\newtcolorbox{acosaserve}[1][A Cosa Serve \& Dove si Usa nella Pratica]{
    enhanced, breakable,
    colback=usebg, colframe=useframe,
    fonttitle=\bfseries\small\color{white},
    title={$\blacktriangleright$ #1}, arc=1.5mm, boxrule=0.8pt,
    left=2.5mm, right=2.5mm, top=2mm, bottom=2mm,
    before skip=4pt, after skip=4pt
}

% Sotto-Box 3: Dimostrazione da Lavagna
\newtcolorbox{dimostrazione}[1][Dimostrazione Passo-Passo (Cosa scrivere alla lavagna)]{
    enhanced, breakable,
    colback=demobg, colframe=demoframe,
    fonttitle=\bfseries\small\color{white},
    title={$\blacktriangleright$ #1}, arc=1.5mm, boxrule=0.8pt,
    left=2.5mm, right=2.5mm, top=2mm, bottom=2mm,
    before skip=4pt, after skip=4pt
}

% Sotto-Box 4: Esposizione Orale
\newtcolorbox{orale}[1][Cosa dire a voce al Prof (Discorso da 30L)]{
    enhanced, breakable,
    colback=oralbg, colframe=oralframe,
    fonttitle=\bfseries\small\color{white},
    title={$\blacktriangleright$ #1}, arc=1.5mm, boxrule=0.8pt,
    left=2.5mm, right=2.5mm, top=2mm, bottom=2mm,
    before skip=4pt, after skip=4pt
}

\newcommand{\PP}{\mathbb{P}}
\newcommand{\E}{\mathbb{E}}
\newcommand{\F}{\mathcal{F}}
\newcommand{\G}{\mathcal{G}}
\newcommand{\1}{\mathbf{1}}

\title{\textbf{\Huge Prontuario di Tutti i Teoremi d'Esame}\\[0.2cm]
\Large Enunciato Formale, Cosa Fa, A Cosa Serve e Dimostrazione da Lavagna\\[0.3cm]
\large \textbf{Corso}: Modelli Probabilistici (518512) --- \textbf{Docente}: Prof. Salvatore Federico}
\author{\textbf{Francesco Castaldi} (\texttt{francesco.castaldi4@studio.unibo.it})\\Corso di Laurea in Informatica --- Università di Bologna}
\date{Anno Accademico 2025/2026 --- Appello di Settembre 2026}

\begin{document}
\maketitle
\tableofcontents
\newpage

% =============================================================================
\chapter{Spazi Finiti, Eventi, Variabili e Valore Atteso}
% =============================================================================

\begin{teoremabox}{Teorema 1.1 --- Proprietà Derivate dagli Assiomi di Kolmogorov}
Sia $(\Omega, \F, \PP)$ uno spazio di probabilità. Per ogni evento $A, B \in \F$:
\begin{enumerate}[leftmargin=*]
    \item $\PP(\emptyset) = 0$ e $\PP(A^c) = 1 - \PP(A)$.
    \item \textbf{Monotonia}: Se $A \subseteq B \implies \PP(A) \le \PP(B)$.
    \item \textbf{Formula di Inclusione-Esclusione}: $\PP(A \cup B) = \PP(A) + \PP(B) - \PP(A \cap B)$.
\end{enumerate}
\end{teoremabox}

\begin{cosafa}
Stabilisce le regole di coerenza logica della misura di probabilità: l'evento impossibile ha probabilità zero, la negazione di un evento ha probabilità complementare a 1, un evento più grande non può avere probabilità inferiore a un suo sottoinsieme, e per sommare due eventi che si sovrappongono bisogna sottrarre l'intersezione per non contarla due volte.
\end{cosafa}

\begin{acosaserve}
È il mattone base di qualsiasi calcolo probabilistico. Serve ogni volta che devi calcolare la probabilità di un'unione di eventi non disgiunti o ricavare la probabilità di un evento raro calcolando il suo contrario ($1 - \PP(A^c)$).
\end{acosaserve}

\begin{dimostrazione}
\begin{enumerate}
    \item $\Omega = A \cup A^c$ con $A \cap A^c = \emptyset \implies \PP(\Omega) = \PP(A) + \PP(A^c) = 1 \implies \PP(A^c) = 1 - \PP(A)$. Ponendo $A=\Omega \implies \PP(\emptyset) = 1 - 1 = 0$.
    \item Se $A \subseteq B$, decomponiamo $B = A \cup (B \setminus A)$ con insiemi disgiunti $\implies \PP(B) = \PP(A) + \PP(B \setminus A) \ge \PP(A)$ poiché $\PP(B \setminus A) \ge 0$.
    \item Scriviamo $A \cup B = A \cup (B \setminus (A \cap B))$ come unione disgiunta $\implies \PP(A \cup B) = \PP(A) + \PP(B \setminus (A \cap B)) = \PP(A) + \PP(B) - \PP(A \cap B)$. $\blacksquare$
\end{enumerate}
\end{dimostrazione}

\begin{orale}
``Professore, queste proprietà fondamentali discendono direttamente dai 3 assiomi di Kolmogorov decomponendo gli insiemi in unioni disgiunte e sfruttando la non-negatività della misura.''
\end{orale}

\vspace{0.4cm}

\begin{teoremabox}{Teorema 1.2 --- Criterio di Misurabilità di Doob}
Siano $X: \Omega \to E$ e $Y: \Omega \to E'$ due variabili aleatorie.\\
La variabile $X$ è $\sigma(Y)$-misurabile $\iff$ esiste una funzione deterministica $f: E' \to E$ tale che $X = f(Y)$.
\end{teoremabox}

\begin{cosafa}
Formalizza matematicamente il concetto di \textbf{dipendenza informativa totale}: dice che $X$ non possiede alcuna incertezza o casualità autonoma rispetto a $Y$. Se conosciamo il valore di $Y$, il valore di $X$ è già perfettamente determinato tramite una funzione $f$.
\end{cosafa}

\begin{acosaserve}
Serve a verificare se una quantità casuale è interamente calcolabile a partire dai dati già osservati. È fondamentale nella teoria dei processi stocastici per definire strategie prevedibili e controlli che usano solo l'informazione storica passata.
\end{acosaserve}

\begin{dimostrazione}
\begin{itemize}
    \item ($\Leftarrow$) Se $X = f(Y)$, per ogni $B \subseteq E$ si ha $\{X \in B\} = \{\omega : f(Y(\omega)) \in B\} = \{Y \in f^{-1}(B)\} \in \sigma(Y)$ per definizione di algebra generata $\sigma(Y)$.
    \item ($\Rightarrow$) Gli atomi di $\sigma(Y)$ sono gli insiemi di livello $\{Y = y\}$. Poiché $X$ è $\sigma(Y)$-misurabile, $X$ deve essere costante su ciascun atomo $\{Y = y\}$. Definiamo la funzione $f(y)$ come l'unico valore assunto da $X$ su tale atomo. Allora $X = f(Y)$. $\blacksquare$
\end{itemize}
\end{dimostrazione}

\newpage

\begin{teoremabox}{Teorema 1.3 --- Legge della Probabilità Totale (Disintegrazione)}
Sia $\{A_1, \dots, A_n\}$ una partizione dello spazio $\Omega$ ($\bigcup A_i = \Omega$ e $A_i \cap A_j = \emptyset$ per $i \ne j$, con $\PP(A_i) > 0$). Per ogni evento $B \in \F$:
\[ \PP(B) = \sum_{i=1}^n \PP(B \cap A_i) = \sum_{i=1}^n \PP(B \mid A_i) \PP(A_i) \]
\end{teoremabox}

\begin{cosafa}
Permette di calcolare la probabilità di un evento complesso $B$ scomponendolo in diversi scenari mutuamente esclusivi $A_i$ (partizione), calcolando la probabilità condizionata di $B$ in ciascuno scenario e pesandola per la probabilità che quello scenario si verifichi.
\end{cosafa}

\begin{acosaserve}
È il cuore della \textbf{First-Step Analysis} nel nostro modello SIR! Per calcolare la probabilità di estinzione o il tempo medio di assorbimento, condizioniamo rispetto a tutti i possibili stati in cui la catena può andare al primo passo $X_1$.
\end{acosaserve}

\begin{dimostrazione}
Poiché $\{A_i\}$ è una partizione di $\Omega$, l'evento $B$ si decompone come:
\[ B = B \cap \Omega = B \cap \left(\bigcup_{i=1}^n A_i\right) = \bigcup_{i=1}^n (B \cap A_i) \]
Poiché gli eventi $(B \cap A_i)$ sono disgiunti a due a due, per l'assioma di additività:
\[ \PP(B) = \sum_{i=1}^n \PP(B \cap A_i) \]
Dalla definizione di probabilità condizionata $\PP(B \cap A_i) = \PP(B \mid A_i)\PP(A_i)$, otteniamo la tesi. $\blacksquare$
\end{dimostrazione}

\vspace{0.4cm}

\begin{teoremabox}{Teorema 1.4 --- Teorema di Bayes (Inversione del Condizionamento)}
Sia $\{A_1, \dots, A_n\}$ una partizione di $\Omega$ e sia $B \in \F$ con $\PP(B) > 0$. Per ogni $k \in \{1, \dots, n\}$:
\[ \PP(A_k \mid B) = \frac{\PP(B \mid A_k)\PP(A_k)}{\PP(B)} = \frac{\PP(B \mid A_k)\PP(A_k)}{\sum_{j=1}^n \PP(B \mid A_j)\PP(A_j)} \]
\end{teoremabox}

\begin{cosafa}
Consente di \textbf{invertire la direzione temporale e causale} del condizionamento: se conosciamo la probabilità di osservare un effetto $B$ data una certa causa $A_k$ ($\PP(B \mid A_k)$), Bayes ci permette di calcolare la probabilità che la causa sia stata $A_k$ avendo osservato l'effetto $B$ ($\PP(A_k \mid B)$).
\end{cosafa}

\begin{acosaserve}
Viene utilizzato nell'inferenza statistica, nella diagnostica medica (calcolo della probabilità di essere malati dato un test positivo) e nell'aggiornamento bayesiano dei parametri di un modello man mano che arrivano nuovi dati.
\end{acosaserve}

\begin{dimostrazione}
Dalla definizione di probabilità condizionata:
\[ \PP(A_k \mid B) = \frac{\PP(A_k \cap B)}{\PP(B)} \]
Riscriviamo il numeratore come $\PP(B \mid A_k)\PP(A_k)$ e il denominatore espandendolo con la Legge della Probabilità Totale $\PP(B) = \sum_{j=1}^n \PP(B \mid A_j)\PP(A_j)$. $\blacksquare$
\end{dimostrazione}

\newpage

\begin{teoremabox}{Teorema 1.5 --- Proprietà della Torre del Valore Atteso Condizionale}
Sia $(\Omega, \F, \PP)$ uno spazio di probabilità e $\G \subseteq \F$ una sub-$\sigma$-algebra. Per ogni $X \in L^1$:
\[ \E\big[\E[X \mid \G]\big] = \E[X] \]
Più in generale, se $\mathcal{H} \subseteq \G$, allora $\E[\E[X \mid \G] \mid \mathcal{H}] = \E[X \mid \mathcal{H}]$.
\end{teoremabox}

\begin{cosafa}
Dice che fare la media ponderata delle migliori previsioni condizionali restituisce esattamente la media globale non condizionata. Inoltre, se componiamo due stime condizionali con livelli di informazione diversi ($\mathcal{H} \subseteq \G$), prevale sempre l'informazione più grossolana $\mathcal{H}$ (Tower Property).
\end{cosafa}

\begin{acosaserve}
È il motore di calcolo principale per le \textbf{Martingale}: permette di dimostrare che la media di una martingala è costante nel tempo ($\E[X_n] = \E[X_0]$) applicando iterativamente la torre sulle filtrazioni $(\F_n)$.
\end{acosaserve}

\begin{dimostrazione}
Dimostriamo il caso di partizione finita $\G = \sigma(A_1, \dots, A_k)$ con $\PP(A_i) > 0$:
\begin{enumerate}
    \item La variabile aleatoria $\E[X \mid \G]$ assume il valore costante $\E[X \mid A_i] = \frac{1}{\PP(A_i)}\sum_{\omega \in A_i} X(\omega)\PP(\{\omega\})$ su ciascun atomo $A_i$.
    \item Calcoliamo il valore atteso della variabile $\E[X \mid \G]$:
    \[ \E\big[\E[X \mid \G]\big] = \sum_{i=1}^k \E[X \mid A_i] \cdot \PP(A_i) = \sum_{i=1}^k \left( \frac{1}{\PP(A_i)} \sum_{\omega \in A_i} X(\omega)\PP(\{\omega\}) \right) \PP(A_i) \]
    \item Semplificando $\PP(A_i)$:
    \[ = \sum_{i=1}^k \sum_{\omega \in A_i} X(\omega)\PP(\{\omega\}) = \sum_{\omega \in \Omega} X(\omega)\PP(\{\omega\}) = \E[X] \quad \blacksquare \]
\end{enumerate}
\end{dimostrazione}

\vspace{0.4cm}

\begin{teoremabox}{Teorema 1.6 --- Estrazione del Noto (Taking out what is known)}
Sia $Y$ una variabile aleatoria $\G$-misurabile e limitata. Allora per ogni $X \in L^1$:
\[ \E[Y X \mid \G] = Y \cdot \E[X \mid \G] \]
\end{teoremabox}

\begin{cosafa}
Dice che all'interno di un'aspettativa condizionata rispetto a $\G$, tutto ciò che è già noto conoscendo $\G$ (cioè $Y$) si comporta come se fosse una costante e può essere ``portato fuori'' dall'operatore di valore atteso.
\end{cosafa}

\begin{acosaserve}
Serve a semplificare drasticamente i calcoli di processi stocastici, strategie di trading $h_n (X_n - X_{n-1})$ e derivazione del compensatore di Doob.
\end{acosaserve}

\begin{dimostrazione}
Sull'atomo $A_i$ di $\G$, poiché $Y$ è $\G$-misurabile, assume un valore costante $y_i$. Allora:
\[ \E[Y X \mid A_i] = \frac{1}{\PP(A_i)} \sum_{\omega \in A_i} Y(\omega)X(\omega)\PP(\{\omega\}) = y_i \left( \frac{1}{\PP(A_i)} \sum_{\omega \in A_i} X(\omega)\PP(\{\omega\}) \right) = y_i \E[X \mid A_i] \]
Ricostruendo la variabile aleatoria su tutti gli atomi: $\E[YX \mid \G] = Y \cdot \E[X \mid \G]$. $\blacksquare$
\end{dimostrazione}

\newpage

% =============================================================================
\chapter{Spazio Discreto Numerabile, Disuguaglianze e Convergenze}
% =============================================================================

\begin{teoremabox}{Teorema 2.1 --- Assenza di Memoria della Distribuzione Geometrica}
Sia $T \sim \text{Geom}(p)$ il tempo di primo successo in prove di Bernoulli indipendenti. Per ogni $n, k \ge 1$:
\[ \PP(T = n + k \mid T > n) = \PP(T = k) \]
\end{teoremabox}

\begin{cosafa}
Dice che la distribuzione geometrica non invecchia mai: se abbiamo già atteso $n$ passi senza successo, la probabilità di dover attendere ulteriori $k$ passi è identica alla probabilità che avevamo di attendere $k$ passi partendo dall'inizio al tempo zero.
\end{cosafa}

\begin{acosaserve}
È la proprietà fondamentale che caratterizza i tempi di permanenza negli stati delle Catene di Markov e i processi di Poisson.
\end{acosaserve}

\begin{dimostrazione}
\begin{enumerate}
    \item La coda della geometrica vale $\PP(T > n) = (1-p)^n$ (significa $n$ fallimenti consecutivi).
    \item L'evento $\{T = n+k\} \cap \{T > n\} = \{T = n+k\}$, con probabilità $(1-p)^{n+k-1}p$.
    \item Applicando la probabilità condizionata:
    \[ \PP(T = n+k \mid T > n) = \frac{\PP(T = n+k)}{\PP(T > n)} = \frac{(1-p)^{n+k-1}p}{(1-p)^n} = (1-p)^{k-1}p = \PP(T = k) \quad \blacksquare \]
\end{enumerate}
\end{dimostrazione}

\vspace{0.4cm}

\begin{teoremabox}{Teorema 2.2 --- Disuguaglianza di Markov}
Sia $X$ una variabile aleatoria non negativa ($X \ge 0$) e sia $a > 0$ una costante. Allora:
\[ \PP(X \ge a) \le \frac{\E[X]}{a} \]
\end{teoremabox}

\begin{cosafa}
Mette un limite superiore invalicabile alla probabilità che una grandezza positiva assuma valori molto grandi, conoscendo unicamente la sua media $\E[X]$.
\end{cosafa}

\begin{acosaserve}
È lo strumento cardine per dimostrare la \textbf{Disuguaglianza di Chebyshev}, il controllo dei rischi e le convergenze stocastiche.
\end{acosaserve}

\begin{dimostrazione}
\begin{enumerate}
    \item Scriviamo il valore atteso spezzando la somma: $\E[X] = \sum_{x < a} x \PP(X=x) + \sum_{x \ge a} x \PP(X=x)$.
    \item Poiché $X \ge 0$, la prima somma è $\ge 0$, dunque $\E[X] \ge \sum_{x \ge a} x \PP(X=x)$.
    \item Per $x \ge a$, minoriamo $x$ con $a$: $\E[X] \ge \sum_{x \ge a} a \PP(X=x) = a \sum_{x \ge a} \PP(X=x) = a \PP(X \ge a)$.
    \item Dividendo per $a > 0$: $\PP(X \ge a) \le \frac{\E[X]}{a}$. $\blacksquare$
\end{enumerate}
\end{dimostrazione}

\vspace{0.4cm}

\begin{teoremabox}{Teorema 2.3 --- Disuguaglianza di Chebyshev}
Sia $X$ una variabile aleatoria con media $m = \E[X]$ e varianza finita $\sigma^2 = \text{Var}(X)$. Per ogni $\epsilon > 0$:
\[ \PP(|X - m| \ge \epsilon) \le \frac{\sigma^2}{\epsilon^2} \]
\end{teoremabox}

\begin{cosafa}
Quantifica la probabilità che una variabile si allontani dalla sua media di più di una distanza $\epsilon$, dimostrando che tale probabilità è controllata dal rapporto tra la varianza $\sigma^2$ e $\epsilon^2$.
\end{cosafa}

\begin{acosaserve}
Serve a dimostrare matematicamente la \textbf{Legge Debole dei Grandi Numeri}: all'aumentare dei campioni la varianza della media campionaria va a zero ($\sigma^2/n \to 0$), costringendo la media campionaria a convergere al valore vero.
\end{acosaserve}

\begin{dimostrazione}
Definiamo la variabile non negativa $Y = (X - m)^2 \ge 0$. L'evento $\{|X - m| \ge \epsilon\}$ equivale a $\{Y \ge \epsilon^2\}$.  
Applicando la disuguaglianza di Markov alla variabile $Y$ con soglia $a = \epsilon^2$:
\[ \PP(|X - m| \ge \epsilon) = \PP(Y \ge \epsilon^2) \le \frac{\E[Y]}{\epsilon^2} = \frac{\E[(X - m)^2]}{\epsilon^2} = \frac{\sigma^2}{\epsilon^2} \quad \blacksquare \]
\end{dimostrazione}

\newpage

\begin{teoremabox}{Teorema 2.4 --- Teorema della Convergenza Dominata di Lebesgue}
Sia $(X_n)_{n \ge 1}$ una successione di variabili aleatorie tali che $X_n \to X$ quasi certamente.\\
Se esiste una variabile aleatoria $G \in L^1$ ($\E[G] < \infty$) tale che $|X_n| \le G$ per ogni $n \ge 1$, allora:
\[ \lim_{n \to \infty} \E[X_n] = \E[X] \]
\end{teoremabox}

\begin{cosafa}
Fornisce la condizione sufficiente definitiva per \textbf{scambiare l'operatore di limite con l'operatore di valore atteso} ($\lim \E[X_n] = \E[\lim X_n]$): basta che le variabili siano dominate da un tetto integrabile $G$.
\end{cosafa}

\begin{acosaserve}
È indispensabile per dimostrare il \textbf{Teorema di Arresto di Doob} sulle Martingale e per giustificare il passaggio al limite per orizzonti temporali infiniti.
\end{acosaserve}

\begin{dimostrazione}
Poiché $|X_n| \le G$, abbiamo $G + X_n \ge 0$ e $G - X_n \ge 0$. Applichiamo il Lemma di Fatou:
\begin{enumerate}
    \item $\E[G + X] = \E[\liminf(G + X_n)] \le \liminf \E[G + X_n] = \E[G] + \liminf \E[X_n] \implies \E[X] \le \liminf \E[X_n]$.
    \item $\E[G - X] = \E[\liminf(G - X_n)] \le \liminf \E[G - X_n] = \E[G] - \limsup \E[X_n] \implies \E[X] \ge \limsup \E[X_n]$.
    \item Dunque $\limsup \E[X_n] \le \E[X] \le \liminf \E[X_n] \implies \lim_{n \to \infty} \E[X_n] = \E[X]$. $\blacksquare$
\end{enumerate}
\end{dimostrazione}

\vspace{0.4cm}

\begin{teoremabox}{Teorema 2.5 --- Legge Debole dei Grandi Numeri (WLLN)}
Sia $(X_n)_{n \ge 1}$ una successione di variabili aleatorie i.i.d. con media $m$ e varianza $\sigma^2 < \infty$.\\
Definita la media campionaria $\bar{S}_n = \frac{1}{n}\sum_{i=1}^n X_i$, per ogni $\epsilon > 0$:
\[ \lim_{n \to \infty} \PP(|\bar{S}_n - m| \ge \epsilon) = 0 \quad (\bar{S}_n \xrightarrow{\PP} m) \]
\end{teoremabox}

\begin{cosafa}
Dimostra che la media aritmetica di $n$ osservazioni casuali indipendenti diventa arbitrariamente vicina alla media teorica della popolazione al crescere del numero di prove $n$.
\end{cosafa}

\begin{acosaserve}
Fornisce la giustificazione teorica di tutte le simulazioni empiriche e degli algoritmi Monte Carlo (come le $M=1000$ traiettorie simulate nel nostro progetto SIR).
\end{acosaserve}

\begin{dimostrazione}
\begin{enumerate}
    \item Calcoliamo media e varianza di $\bar{S}_n$: $\E[\bar{S}_n] = \frac{1}{n}\sum \E[X_i] = m$ e $\text{Var}(\bar{S}_n) = \frac{1}{n^2}\sum \text{Var}(X_i) = \frac{\sigma^2}{n}$.
    \item Applichiamo la disuguaglianza di Chebyshev a $\bar{S}_n$: $\PP(|\bar{S}_n - m| \ge \epsilon) \le \frac{\text{Var}(\bar{S}_n)}{\epsilon^2} = \frac{\sigma^2}{n \epsilon^2}$.
    \item Passando al limite per $n \to \infty$: $\lim_{n \to \infty} \frac{\sigma^2}{n \epsilon^2} = 0$. $\blacksquare$
\end{enumerate}
\end{dimostrazione}

\newpage

% =============================================================================
\chapter{Processi Stocastici, Filtrazioni e Martingale}
% =============================================================================

\begin{teoremabox}{Teorema 3.1 --- Proprietà di Chiusura dei Tempi di Arresto}
Siano $\tau, \sigma$ due tempi di arresto rispetto a $(\F_n)_{n \ge 0}$. Allora:
\begin{enumerate}
    \item $\tau \wedge \sigma = \min(\tau, \sigma)$ è un tempo di arresto.
    \item $\tau \vee \sigma = \max(\tau, \sigma)$ è un tempo di arresto.
    \item $\tau + k$ (con $k \in \mathbb{N}$) è un tempo di arresto.
\end{enumerate}
\end{teoremabox}

\begin{cosafa}
Garantisce che combinando più regole di arresto legali (il primo evento che accade, l'ultimo che accade, o ritardare l'arresto di $k$ giorni), la regola risultante rimane ancora un tempo di arresto valido senza usare informazioni future.
\end{cosafa}

\begin{acosaserve}
Permette di costruire tempi di arresto troncati a un orizzonte finito $\tau \wedge n = \min(\tau, n)$, indispensabili per le dimostrazioni sulle martingale limitate.
\end{acosaserve}

\begin{dimostrazione}
\begin{enumerate}
    \item $\{\tau \wedge \sigma \le n\} = \{\tau \le n\} \cup \{\sigma \le n\} \in \F_n$ poiché unione di elementi di $\F_n$.
    \item $\{\tau \vee \sigma \le n\} = \{\tau \le n\} \cap \{\sigma \le n\} \in \F_n$ poiché intersezione di elementi di $\F_n$.
    \item $\{\tau + k \le n\} = \{\tau \le n - k\} \in \F_{n-k} \subseteq \F_n$. $\blacksquare$
\end{enumerate}
\end{dimostrazione}

\vspace{0.4cm}

\begin{teoremabox}{Teorema 4.1 --- Conservazione della Media di una Martingala}
Sia $X = (X_n)_{n \ge 0}$ una martingala rispetto a $(\F_n)$. Allora per ogni $n \ge 0$:
\[ \E[X_n] = \E[X_0] \]
\end{teoremabox}

\begin{cosafa}
Dice che in un processo a gioco equo, la ricchezza media attesa rimane costante a qualunque istante futuro $n$, coincidente con il capitale iniziale.
\end{cosafa}

\begin{acosaserve}
È la proprietà cardine utilizzata per calcolare le probabilità di assorbimento e i valori finali dei giochi equi.
\end{acosaserve}

\begin{dimostrazione}
Per induzione su $n$: per $n=0$ è banale. Supponiamo $\E[X_n] = \E[X_0]$.  
Applicando la Tower Property e la definizione di martingala $\E[X_{n+1} \mid \F_n] = X_n$:
\[ \E[X_{n+1}] = \E\big[\E[X_{n+1} \mid \F_n]\big] = \E[X_n] = \E[X_0] \quad \blacksquare \]
\end{dimostrazione}

\vspace{0.4cm}

\begin{teoremabox}{Teorema 4.2 --- Teorema della Martingala Arrestata}
Se $X = (X_n)_{n \ge 0}$ è una martingala e $\tau$ è un tempo di arresto, allora il processo arrestato $X^\tau = (X_{n \wedge \tau})_{n \ge 0}$ è anch'esso una \textbf{martingala}.
\end{teoremabox}

\begin{cosafa}
Dimostra che congelare un gioco equo al momento in cui scatta una regola di stop $\tau$ produce un nuovo processo che continua a essere un gioco equo per tutti i tempi successivi.
\end{cosafa}

\begin{acosaserve}
È il lemma ponte per dimostrare il Teorema di Arresto Opzionale di Doob.
\end{acosaserve}

\begin{dimostrazione}
Scriviamo l'incremento di $X^\tau$ come $X_{n+1 \wedge \tau} - X_{n \wedge \tau} = (X_{n+1} - X_n) \1_{\{\tau > n\}}$.  
Poiché $\{\tau > n\} = \{\tau \le n\}^c \in \F_n$, la variabile indicatrice è nota al tempo $n$ ($\F_n$-misurabile).  
Calcolando il valore atteso condizionato:
\[ \E[X_{n+1 \wedge \tau} - X_{n \wedge \tau} \mid \F_n] = \1_{\{\tau > n\}} \E[X_{n+1} - X_n \mid \F_n] = \1_{\{\tau > n\}} \cdot 0 = 0 \]
Dunque $\E[X_{n+1 \wedge \tau} \mid \F_n] = X_{n \wedge \tau}$. $\blacksquare$
\end{dimostrazione}

\newpage

\begin{teoremabox}{Teorema 4.3 --- Teorema di Arresto Opzionale di Doob}
Sia $X$ una martingala e $\tau$ un tempo di arresto. Se vale una delle seguenti condizioni:
\begin{enumerate}
    \item $\tau$ è limitato quasi certamente ($\tau \le K < \infty$),
    \item $X$ è uniformemente limitato ($|X_{n \wedge \tau}| \le M$) e $\PP(\tau < \infty) = 1$,
    \item $\E[\tau] < \infty$ e gli incrementi sono limitati ($|X_{n+1} - X_n| \le C$),
\end{enumerate}
allora il valore atteso al tempo di arresto si conserva:
\[ \E[X_\tau] = \E[X_0] \]
\end{teoremabox}

\begin{cosafa}
Stabilisce che nessuna regola di arresto (stopping time) può permettere a un giocatore di andarsene con un guadagno medio positivo da un gioco equo: la media della ricchezza all'arresto $\E[X_\tau]$ è identica alla ricchezza iniziale $\E[X_0]$.
\end{cosafa}

\begin{acosaserve}
È lo strumento principe per risolvere problemi di probabilità di uscita da intervalli (es. Rovina del Giocatore) e calcolare tempi medi di primo raggiungimento.
\end{acosaserve}

\begin{dimostrazione}
Per il Teorema della Martingala Arrestata, $\E[X_{n \wedge \tau}] = \E[X_0]$ per ogni $n$.  
Facendo tendere $n \to \infty$, $X_{n \wedge \tau} \to X_\tau$ quasi certamente poiché $\PP(\tau < \infty)=1$.  
Sotto l'ipotesi di limitatezza $|X_{n \wedge \tau}| \le M$, applichiamo il Teorema della Convergenza Dominata di Lebesgue:
\[ \E[X_\tau] = \E\left[\lim_{n \to \infty} X_{n \wedge \tau}\right] = \lim_{n \to \infty} \E[X_{n \wedge \tau}] = \E[X_0] \quad \blacksquare \]
\end{dimostrazione}

\newpage

% =============================================================================
\chapter{Passeggiate Aleatorie e Rovina del Giocatore}
% =============================================================================

\begin{teoremabox}{Teorema 5.1 --- Lemmi di Borel-Cantelli}
Sia $(\Omega, \F, \PP)$ uno spazio di probabilità e sia $(A_n)_{n \ge 1}$ una successione di eventi.
\begin{enumerate}
    \item \textbf{Primo Lemma}: Se $\sum_{n=1}^\infty \PP(A_n) < \infty \implies \PP(A_n \text{ infinitamente spesso}) = 0$.
    \item \textbf{Secondo Lemma}: Se $(A_n)$ sono indipendenti e $\sum_{n=1}^\infty \PP(A_n) = \infty \implies \PP(A_n \text{ infinitamente spesso}) = 1$.
\end{enumerate}
\end{teoremabox}

\begin{cosafa}
Fornisce un criterio definitivo basato sulla convergenza o divergenza della serie delle probabilità per stabilire se una sequenza infinita di eventi si verificherà infinite volte nel tempo o se smetterà di verificarsi dopo un tempo finito.
\end{cosafa}

\begin{acosaserve}
Serve a studiare la ricorrenza delle passeggiate aleatorie, la convergenza quasi certa e le proprietà di coda dei processi stocastici.
\end{acosaserve}

\begin{dimostrazione}
Dimostrazione del Primo Lemma:
\begin{enumerate}
    \item L'evento $\{A_n \text{ i.o.}\} = \limsup A_n = \bigcap_{k=1}^\infty \bigcup_{n=k}^\infty A_n$.
    \item Per monotonia e $\sigma$-subadditività: $\PP(\{A_n \text{ i.o.}\}) \le \PP(\bigcup_{n=k}^\infty A_n) \le \sum_{n=k}^\infty \PP(A_n)$.
    \item Poiché la serie converge, la coda tende a zero: $\lim_{k \to \infty} \sum_{n=k}^\infty \PP(A_n) = 0$. Dunque $\PP(A_n \text{ i.o.}) = 0$. $\blacksquare$
\end{enumerate}
\end{dimostrazione}

\vspace{0.4cm}

\begin{teoremabox}{Teorema 5.2 --- La Rovina del Giocatore (Gambler's Ruin)}
In una passeggiata aleatoria simmetrica su $\mathbb{Z}$ ($p=1/2$) con barriere assorbenti a $-a$ (capitale di A) e $+b$ (capitale del banco B):
\[ \PP(\text{Rovina di A}) = \frac{b}{a+b}, \qquad \PP(\text{Rovina di B}) = \frac{a}{a+b} \]
Se il banco ha capitale infinito ($b \to \infty$), $\lim_{b \to \infty} \PP(\text{Rovina di A}) = 1$.
\end{teoremabox}

\begin{cosafa}
Calcola la probabilità esatta che un giocatore vada in bancarotta prima di raddoppiare o vincere il capitale del banco in un gioco d'azzardo equo.
\end{cosafa}

\begin{acosaserve}
È l'applicazione pratica più celebre del Teorema di Arresto Opzionale di Doob e dimostra matematicamente perché non si può vincere contro un casinò con capitale illimitato.
\end{acosaserve}

\begin{dimostrazione}
\begin{enumerate}
    \item Sia $X_n$ la passeggiata simmetrica con $X_0 = 0$. $X_n$ è una martingala limitata nell'intervallo $[-a, b]$.
    \item Il tempo di arresto $\tau = \inf\{n : X_n = -a \text{ oppure } X_n = b\}$ soddisfa $\PP(\tau < \infty) = 1$.
    \item Per il Teorema di Arresto di Doob: $\E[X_\tau] = \E[X_0] = 0$.
    \item Esplicitando il valore atteso: $\E[X_\tau] = (-a)\PP(X_\tau = -a) + b \PP(X_\tau = b) = 0$.
    \item Ponendo $p_A = \PP(X_\tau = -a)$ e sapendo che $p_B = 1 - p_A$:
    \[ -a p_A + b(1 - p_A) = 0 \implies (a+b)p_A = b \implies p_A = \frac{b}{a+b} \quad \blacksquare \]
\end{enumerate}
\end{dimostrazione}

\newpage

% =============================================================================
\chapter{Catene di Markov su Spazio Discreto}
% =============================================================================

\begin{teoremabox}{Teorema 6.1 --- Equazioni di Chapman-Kolmogorov}
Sia $X = (X_n)_{n \ge 0}$ una Catena di Markov omogenea con matrice di transizione $P$. Per ogni $n, m \ge 0$:
\[ P^{(n+m)} = P^n \cdot P^m = P^{n+m}, \qquad p_{ij}^{(n+m)} = \sum_{k \in \mathcal{S}} p_{ik}^{(n)} p_{kj}^{(m)} \]
\end{teoremabox}

\begin{cosafa}
Descrive come calcolare le probabilità di transizione a lungo termine ($n+m$ passi) sommando su tutti i possibili stati intermedi $k$ in cui il sistema poteva trovarsi al passo $n$.
\end{cosafa}

\begin{acosaserve}
Permette di calcolare la distribuzione dello stato della catena a qualsiasi istante futuro tramite la semplice potenza di matrice $\mu_n = \mu_0 P^n$.
\end{acosaserve}

\begin{dimostrazione}
\begin{enumerate}
    \item Per definizione: $p_{ij}^{(n+m)} = \PP(X_{n+m} = j \mid X_0 = i)$.
    \item Disintegriamo rispetto allo stato intermedio $X_n$: $\PP(X_{n+m} = j \mid X_0 = i) = \sum_{k} \PP(X_{n+m} = j, X_n = k \mid X_0 = i)$.
    \item Usiamo la probabilità condizionata: $= \sum_{k} \PP(X_{n+m} = j \mid X_n = k, X_0 = i) \cdot \PP(X_n = k \mid X_0 = i)$.
    \item Per la \textbf{Proprietà di Markov}, il futuro a tempo $n+m$ dato il presente a tempo $n$ non dipende dal passato $X_0 = i$: $\PP(X_{n+m} = j \mid X_n = k, X_0 = i) = p_{kj}^{(m)}$.
    \item Sostituendo: $p_{ij}^{(n+m)} = \sum_{k} p_{ik}^{(n)} p_{kj}^{(m)} = (P^n P^m)_{ij}$. $\blacksquare$
\end{enumerate}
\end{dimostrazione}

\vspace{0.4cm}

\begin{teoremabox}{Teorema 6.2 --- Assorbimento Quasi Certo in Catene Finite}
Sia $X$ una Catena di Markov a stati finiti con partizione $E = \mathcal{T} \cup \mathcal{A}$. Se da ogni stato transitorio $i \in \mathcal{T}$ è possibile raggiungere $\mathcal{A}$ in un numero finito di passi, allora:
\[ \PP(\tau < \infty) = 1 \qquad \text{dove } \tau = \inf\{n \ge 0 : X_n \in \mathcal{A}\} \]
\end{teoremabox}

\begin{cosafa}
Garantisce che in una catena finita con stati assorbenti accessibili, il processo non rimarrà mai intrappolato per sempre negli stati transitori: l'assorbimento avviene con certezza assoluta (probabilità 100\%).
\end{cosafa}

\begin{acosaserve}
È il teorema fondamentale su cui si fonda la conclusione del nostro \textbf{Modello SIR}: dimostra che l'epidemia deve estinguersi prima o poi quasi certamente ($\PP(\tau_{\text{estinzione}} < \infty) = 1$).
\end{acosaserve}

\begin{dimostrazione}
\begin{enumerate}
    \item Sia $m = |\mathcal{T}|$ il numero di stati transitori. Per ogni $i \in \mathcal{T}$ esiste un cammino di lunghezza $\le m$ verso $\mathcal{A}$.
    \item Sia $\epsilon = \min_{i \in \mathcal{T}} \PP_i(\tau \le m) > 0$ la minima probabilità di assorbimento in $m$ passi.
    \item La probabilità di rimanere in $\mathcal{T}$ dopo $k \cdot m$ passi decade geometricamente: $\PP_i(\tau > k \cdot m) \le (1 - \epsilon)^k$.
    \item Poiché $1 - \epsilon < 1$, passando al limite per $k \to \infty$: $\lim_{k \to \infty} (1-\epsilon)^k = 0$.
    \item Dunque $\PP_i(\tau = \infty) = 0 \implies \PP_i(\tau < \infty) = 1$. $\blacksquare$
\end{enumerate}
\end{dimostrazione}

\newpage

\begin{teoremabox}{Teorema 6.3 --- Invertibilità di $(I-Q)$ e Serie di Neumann}
Sia $Q$ la sottomatrice delle transizioni tra stati transitori. Allora il raggio spettrale $\rho(Q) < 1$, $\lim_{k \to \infty} Q^k = 0$, la matrice $(I-Q)$ è invertibile e:
\[ (I-Q)^{-1} = \sum_{k=0}^\infty Q^k \]
\end{teoremabox}

\begin{cosafa}
Dimostra che la serie geometrica matriciale $\sum Q^k$ converge alla \textbf{Matrice Fondamentale} $(I-Q)^{-1}$, dove l'elemento $(i,j)$ conta il numero medio di visite allo stato transitorio $j$ partendo da $i$.
\end{cosafa}

\begin{acosaserve}
È la base algebrica per calcolare tempi medi di assorbimento e probabilità finali di uscita nel modello SIR.
\end{acosaserve}

\begin{dimostrazione}
\begin{enumerate}
    \item Dall'assorbimento quasi certo, $\lim_{k \to \infty} Q^k = 0$.
    \item Consideriamo la somma telescopica finita: $(I - Q) \left( \sum_{k=0}^{n-1} Q^k \right) = I - Q^n$.
    \item Passando al limite per $n \to \infty$, poiché $Q^n \to 0$, il membro di destra tende alla matrice identità $I$.
    \item Dunque la serie $\sum_{k=0}^\infty Q^k$ converge ed è l'inversa di $(I-Q)$. $\blacksquare$
\end{enumerate}
\end{dimostrazione}

\vspace{0.4cm}

\begin{teoremabox}{Teorema 6.4 --- Tempo Medio di Assorbimento via First-Step Analysis}
Il vettore $\mathbf{t} = (t_i)_{i \in \mathcal{T}}$ dei tempi medi di assorbimento partendo da ciascuno stato transitorio soddisfa il sistema lineare:
\[ (I - Q)\mathbf{t} = \mathbf{1} \iff \mathbf{t} = (I - Q)^{-1}\mathbf{1} \]
\end{teoremabox}

\begin{cosafa}
Fornisce la formula esatta per calcolare la durata media attesa di un processo prima di essere assorbito (es. la durata media dell'epidemia SIR).
\end{cosafa}

\begin{acosaserve}
Nel nostro progetto SIR, per $N=3$ abbiamo usato questa formula per trovare il vettore dei tempi medi di tutti i 6 stati transitori invertendo la matrice $6 \times 6$.
\end{acosaserve}

\begin{dimostrazione}
\begin{enumerate}
    \item Sia $t_i = \E[\tau \mid X_0 = i]$. Condizioniamo rispetto al primo passo $X_1$:
    \[ t_i = 1 + \sum_{j \in E} P_{ij} \E[\tau \mid X_1 = j] \]
    \item Separiamo gli stati transitori $\mathcal{T}$ e gli stati assorbenti $\mathcal{A}$:
    \[ t_i = 1 + \sum_{j \in \mathcal{T}} Q_{ij} t_j + \sum_{a \in \mathcal{A}} R_{ia} \cdot 0 = 1 + \sum_{j \in \mathcal{T}} Q_{ij} t_j \]
    \item In forma matriciale: $\mathbf{t} = \mathbf{1} + Q \mathbf{t} \iff (I - Q)\mathbf{t} = \mathbf{1}$.
    \item Moltiplicando a sinistra per $(I-Q)^{-1}$: $\mathbf{t} = (I - Q)^{-1}\mathbf{1}$. $\blacksquare$
\end{enumerate}
\end{dimostrazione}

\vspace{0.4cm}

\begin{teoremabox}{Teorema 6.5 --- Matrice delle Probabilità di Assorbimento}
La matrice $B = (b_{ia})_{i \in \mathcal{T}, a \in \mathcal{A}}$, dove $b_{ia} = \PP(X_\tau = a \mid X_0 = i)$, è data da:
\[ B = (I - Q)^{-1} R \]
\end{teoremabox}

\begin{cosafa}
Calcola la probabilità esatta di finire in ciascuno specifico stato assorbente $a$ partendo dallo stato iniziale $i$.
\end{cosafa}

\begin{acosaserve}
Nel modello SIR permette di calcolare la probabilità analitica esatta dell'attacco finale $R_\infty$ (es. finire in $(2,0,1)$ vs $(0,0,3)$ per $N=3$).
\end{acosaserve}

\begin{dimostrazione}
\begin{enumerate}
    \item Condizioniamo al primo passo $X_1$: $b_{ia} = R_{ia} + \sum_{j \in \mathcal{T}} Q_{ij} b_{ja}$.
    \item In forma matriciale: $B = R + QB \iff (I - Q)B = R$.
    \item Moltiplicando a sinistra per $(I-Q)^{-1}$: $B = (I - Q)^{-1} R$. $\blacksquare$
\end{enumerate}
\end{dimostrazione}

\vspace{0.4cm}

\begin{teoremabox}{Teorema 6.6 --- Teorema Ergodico per Catene di Markov}
Sia $X$ una Catena di Markov irriducibile e positivo-ricorrente su $\mathcal{S}$ con misura di probabilità invariante $\pi$ ($\pi P = \pi$). Allora per ogni funzione limitata $f: \mathcal{S} \to \mathbb{R}$:
\[ \lim_{n \to \infty} \frac{1}{n} \sum_{k=1}^n f(X_k) = \sum_{x \in \mathcal{S}} f(x) \pi(x) \quad \text{quasi certamente} \]
\end{teoremabox}

\begin{cosafa}
Stabilisce l'equivalenza fondamentale tra media temporale lungo una singola traiettoria e media spaziale pesata con la distribuzione stazionaria $\pi$.
\end{cosafa}

\begin{acosaserve}
È il fondamento teorico dei metodi MCMC (Markov Chain Monte Carlo) e degli algoritmi di campionamento stocastico.
\end{acosaserve}

\end{document}
